\documentclass[12pt,letterpaper]{article}

% ── Encoding ────────────────────────────────────────────────────────────────
\usepackage[utf8]{inputenc}
\usepackage[T1]{fontenc}

% ── Packages ────────────────────────────────────────────────────────────────
\usepackage[margin=1in]{geometry}
\usepackage{hyperref}
\usepackage{xcolor}
\usepackage{titlesec}
\usepackage{enumitem}
\usepackage{booktabs}
\usepackage{longtable}
\usepackage{array}
\usepackage{fancyhdr}
\usepackage{graphicx}
\usepackage{listings}
\usepackage{courier}
\usepackage{mdframed}
\usepackage{tcolorbox}
\usepackage{amsmath}
\usepackage{amssymb}
\usepackage{microtype}
\usepackage{parskip}
\usepackage{float}

% ── Colors ──────────────────────────────────────────────────────────────────
\definecolor{njitred}{RGB}{204,0,0}
\definecolor{darkgray}{RGB}{51,51,51}
\definecolor{lightgray}{RGB}{245,245,245}
\definecolor{codegray}{RGB}{240,240,240}
\definecolor{commentgreen}{RGB}{0,128,0}
\definecolor{keywordblue}{RGB}{0,0,180}
\definecolor{stringbrown}{RGB}{160,82,45}

% ── Hyperref setup ──────────────────────────────────────────────────────────
\hypersetup{
  colorlinks=true,
  linkcolor=njitred,
  urlcolor=njitred,
  citecolor=njitred,
  pdftitle={GSA Gateway Technical Reference},
  pdfauthor={Mohammad Dindoost},
}

% ── Code listing style ───────────────────────────────────────────────────────
\lstset{
  basicstyle=\ttfamily\footnotesize,
  backgroundcolor=\color{codegray},
  frame=single,
  rulecolor=\color{darkgray!40},
  breaklines=true,
  breakatwhitespace=true,
  tabsize=4,
  showstringspaces=false,
  keywordstyle=\color{keywordblue}\bfseries,
  commentstyle=\color{commentgreen}\itshape,
  stringstyle=\color{stringbrown},
  numbers=left,
  numberstyle=\tiny\color{darkgray},
  numbersep=5pt,
  captionpos=b,
  extendedchars=true,
  inputencoding=utf8,
  literate=
    {->}{{$\rightarrow$}}{2}
    {->\ }{{$\rightarrow$\ }}{3}
    {-->}{{$\longrightarrow$}}{3}
    {->|}{{$\rightarrow$|}}{3}
    {<-}{{$\leftarrow$}}{2}
    {=>}{{$\Rightarrow$}}{2}
    {~}{{$\sim$}}{1}
    {<=}{{$\leq$}}{2}
    {>=}{{$\geq$}}{2}
    {!=}{{$\neq$}}{2}
    {|}{{$|$}}{1}
    {``}{{\textquoteleft\textquoteleft}}{2}
    {''}{{"}}{2}
    {"}{{"}}{1}
    {"}{{"}}{1}
    {'}{{'}}{1}
    {'}{{'}}{1}
    {--}{{--}}{2}
    {---}{{---}}{3}
    {*}{{\textasteriskcentered}}{1}
    {euro}{{\euro}}{4}
    {<==>}{{$\Longleftrightarrow$}}{4}
    {<=>}{{$\Leftrightarrow$}}{3}
    ,
}

% ── Section title colors ─────────────────────────────────────────────────────
\titleformat{\section}
  {\Large\bfseries\color{njitred}}{\thesection}{1em}{}[\color{njitred}\titlerule]
\titleformat{\subsection}
  {\large\bfseries\color{darkgray}}{\thesubsection}{1em}{}
\titleformat{\subsubsection}
  {\normalsize\bfseries\color{darkgray}}{\thesubsubsection}{1em}{}

% ── Header / Footer ─────────────────────────────────────────────────────────
\pagestyle{fancy}
\fancyhf{}
\fancyhead[L]{\small\textcolor{njitred}{\textbf{GSA Gateway}} — Technical Reference}
\fancyhead[R]{\small NJIT Graduate Student Association}
\fancyfoot[C]{\thepage}
\renewcommand{\headrulewidth}{0.4pt}
\renewcommand{\headrule}{\hbox to\headwidth{\color{njitred}\leaders\hrule height \headrulewidth\hfill}}

% ── Highlight box ────────────────────────────────────────────────────────────
\tcbuselibrary{skins,breakable}
\newtcolorbox{infobox}[1][]{
  colback=lightgray,
  colframe=njitred,
  fonttitle=\bfseries,
  title=#1,
  breakable,
}
\newtcolorbox{warningbox}[1][]{
  colback=yellow!10,
  colframe=orange!80!black,
  fonttitle=\bfseries,
  title=#1,
  breakable,
}

% ─────────────────────────────────────────────────────────────────────────────
\begin{document}

% ── Title Page ───────────────────────────────────────────────────────────────
\begin{titlepage}
  \centering
  \vspace*{2cm}
  {\Huge\bfseries\color{njitred} GSA Gateway\\[0.5em]}
  {\Large\color{darkgray} Technical Reference for New Contributors}\\[2em]
  \rule{0.6\textwidth}{0.4pt}\\[1.5em]
  {\large New Jersey Institute of Technology\\
   Graduate Student Association\\[1em]
   Discord Bot + Static Website}\\[2em]
  {\normalsize
    \textbf{Author:} Mohammad Dindoost\\
    \textbf{Role:} VP Academic Affairs, Ph.D.\ Candidate in Computer Science\\
    \textbf{Contact:} \href{mailto:md724@njit.edu}{md724@njit.edu}\\[1em]
    \textbf{Repository:} \texttt{gsa-gateway}\\
    \textbf{Branch:} \texttt{main}\\[1em]
    \textbf{Document Version:} June 2026
  }\\[3em]
  \rule{0.6\textwidth}{0.4pt}\\[1em]
  {\small This document is intended for new contributors, student developers, and
   researchers who want to understand how GSA Gateway is built, why specific
   technology choices were made, and how to extend the system.}
\end{titlepage}

\tableofcontents
\newpage

% =============================================================================
\section{Project Overview}
% =============================================================================

GSA Gateway is the official digital infrastructure of the NJIT Graduate Student
Association. It consists of two independently deployable components:

\begin{enumerate}[leftmargin=2em]
  \item \textbf{Discord Bot} — A Python bot that answers student questions using a
        full Retrieval-Augmented Generation (RAG) pipeline, manages slash commands for
        events, resources, feedback, and travel funding, and runs automated daily
        digests.
  \item \textbf{Static Website} — A pure HTML/CSS/JS website hosted on GitHub Pages
        that surfaces the same information (events, officers, resources) to students
        without Discord.
\end{enumerate}

\subsection{What Problems Does It Solve?}

Before GSA Gateway, students had no single place to answer questions like:
\begin{itemize}[leftmargin=2em]
  \item ``How do I apply for a travel award and what is the deadline?''
  \item ``Who is the GSA VP of Finances and how do I reach them?''
  \item ``What events are happening this month?''
  \item ``What is the reimbursement cap for international travel?''
\end{itemize}

Officers answered these repetitively via email and Discord messages. GSA Gateway
automates this by grounding every answer in official documents and routing
questions intelligently.

\subsection{Technology Stack at a Glance}

\begin{center}
\begin{tabular}{lll}
\toprule
\textbf{Component} & \textbf{Technology} & \textbf{Why} \\
\midrule
Bot framework & discord.py 2.x & Mature, async-native, slash command support \\
Language & Python 3.11+ & Async ecosystem, ML libraries \\
LLM (generation) & Ollama + llama3.1:8b & Local, no API cost, privacy-safe \\
Embedding model & nomic-embed-text & Open source, strong semantic quality \\
Vector database & ChromaDB & Persistent, zero-server, Python-native \\
Relational database & SQLite & Simple, file-based, no infrastructure \\
Fuzzy search & rapidfuzz & Fast approximate string matching \\
Tokenizer & tiktoken (cl100k\_base) & Accurate token counting for chunking \\
Website & HTML/CSS/JS & No framework, GitHub Pages compatible \\
\bottomrule
\end{tabular}
\end{center}

% =============================================================================
\section{Repository Structure}
% =============================================================================

\begin{lstlisting}[language=bash, caption={Top-level directory layout}]
gsa-gateway/
+-- bot/                    # All Python bot code
|   +-- main.py             # Entry point -- GSABot class, startup sequence
|   +-- config.py           # Reads .env -> typed Config dataclass
|   +-- commands/           # Discord slash commands (cogs)
|   +-- services/           # Business logic, RAG pipeline, database
|   \-- data/               # Knowledge base files (Markdown + YAML)
+-- website/                # Static site
|   +-- *.html              # One file per page
|   +-- assets/             # style.css, app.js
|   \-- data/events.json    # Auto-generated from events.yml
+-- scripts/                # Utility scripts (setup, index build, export)
+-- docs/                   # Architecture and deployment documentation
+-- chroma_db/              # ChromaDB persistent store (auto-created)
+-- gsa_gateway.db          # SQLite database (auto-created)
+-- gsa_gateway.log         # Runtime log (auto-created)
\-- .env                    # Local secrets (never committed)
\end{lstlisting}

\subsection{The \texttt{bot/commands/} Directory}

Each file here is a Discord \emph{cog} — a self-contained class that registers
one or more slash commands. Adding a new command means creating a new file and
listing it in \texttt{EXTENSIONS} in \texttt{main.py}. No other wiring needed.

\begin{center}
\begin{longtable}{lp{9cm}}
\toprule
\textbf{File} & \textbf{Responsibility} \\
\midrule
\texttt{ask.py} & \texttt{/ask} — fuzzy-searches knowledge base, optional Ollama LLM answer \\
\texttt{chat.py} & \texttt{on\_message} handler — free-form chat in \#ask-gsa and DMs; full RAG pipeline \\
\texttt{events.py} & \texttt{/events} (list) and \texttt{/event [name]} (detail) \\
\texttt{initiative.py} & \texttt{/initiative} — Discord Modal collecting 5 fields, stored to SQLite \\
\texttt{feedback.py} & \texttt{/feedback} — anonymous message stored to SQLite \\
\texttt{resources.py} & \texttt{/resources [category]} — lists YAML-sourced resources \\
\texttt{contact.py} & \texttt{/contact [role]} — directory lookup from \texttt{contacts.yml} \\
\texttt{help\_cmd.py} & \texttt{/help} — ephemeral embed with full command reference \\
\texttt{admin.py} & \texttt{/admin\_summary}, \texttt{/admin\_export}, \texttt{/admin\_stats}, \texttt{/admin\_announce}, \texttt{/admin\_add\_event}, \texttt{/admin\_rebuild\_index} (all ephemeral, admin role gated) \\
\texttt{qrcode\_cmd.py} & \texttt{/qrcode} — generates GSA-branded QR code with color options \\
\bottomrule
\end{longtable}
\end{center}

\subsection{The \texttt{bot/services/} Directory}

Services are the business logic layer. They have no knowledge of Discord; they
receive plain Python objects and return plain Python objects. This separation
makes them independently testable.

\begin{center}
\begin{longtable}{lp{9cm}}
\toprule
\textbf{File} & \textbf{Responsibility} \\
\midrule
\texttt{chunker.py} & Splits all KB files into $\leq$350-token chunks for vector storage \\
\texttt{embedder.py} & Converts text to vectors using Ollama's \texttt{nomic-embed-text} \\
\texttt{vector\_store.py} & Wraps ChromaDB — stores chunks, queries by cosine similarity \\
\texttt{bm25\_index.py} & BM25Okapi lexical index over all chunks — exact/sparse term retrieval \\
\texttt{retriever.py} & Hybrid: vector + BM25 in parallel, RRF fusion, rerank to top-7, expand siblings \\
\texttt{ollama\_client.py} & Sends prompts to Ollama LLM, builds system/user prompt from chunks \\
\texttt{conversation.py} & Per-user in-memory session history (60 min timeout, 5-turn window) \\
\texttt{intent\_detector.py} & Classifies messages — greeting, food, thanks, help, question, statement \\
\texttt{knowledge\_base.py} & Loads all \texttt{.md} and \texttt{.yml} data files into memory at startup \\
\texttt{search.py} & rapidfuzz fuzzy search over KB entries (threshold: 60.0) \\
\texttt{database.py} & All SQLite CRUD operations, privacy hashing, statistics \\
\texttt{moderation.py} & \texttt{RateLimiter} (5 calls/60s), channel allowlist, admin role check \\
\texttt{summaries.py} & \texttt{weekly\_summary()} — generates markdown summary from SQLite data \\
\texttt{mathcafe.py} & Posts a daily math fact to \#gsa-mathcafe; tracks Maryam Mirzakhani birthday \\
\texttt{scheduler.py} & Daily digest cog — posts announcements on schedule \\
\bottomrule
\end{longtable}
\end{center}

\subsection{The \texttt{bot/data/} Directory}

This is the knowledge base. Every fact the bot can answer comes from one of
these files. Editing a Markdown or YAML file here changes what the bot knows.
After editing, you must rebuild the vector index.

\begin{center}
\begin{longtable}{lll}
\toprule
\textbf{File} & \textbf{Format} & \textbf{Content} \\
\midrule
\texttt{gsa\_faq.md} & Markdown FAQ & 12+ Q\&A pairs about GSA services \\
\texttt{mmi\_workshop.md} & Markdown FAQ & Full MMI 2025 and 2026 knowledge base \\
\texttt{bot\_features.md} & Markdown FAQ & Bot command reference for \#ask-gsa \\
\texttt{gsa\_constitution.md} & Markdown policy & GSA constitution and bylaws \\
\texttt{travel\_award.md} & Markdown policy & Travel award process, tiers, Chrome River \\
\texttt{club\_finance.md} & Markdown policy & Club financial bylaws \\
\texttt{rules.md} & Markdown policy & Community guidelines \\
\texttt{events.yml} & YAML & 5+ events with dates, locations, descriptions \\
\texttt{contacts.yml} & YAML & 9 officers and campus offices with emails \\
\texttt{resources.yml} & YAML & 8 categories $\times$ 3–4 resources each \\
\bottomrule
\end{longtable}
\end{center}

\begin{infobox}[Markdown FAQ Format]
All \texttt{.md} FAQ files must follow this exact format so the chunker can
parse them correctly:
\begin{lstlisting}
## Q: Your question here?
**A:** Your answer here (can span multiple lines).
\end{lstlisting}
\end{infobox}

% =============================================================================
\section{System Architecture}
% =============================================================================

\subsection{Startup Sequence}

When the bot starts, \texttt{GSABot.setup\_hook()} runs before Discord marks the
bot as online. This ensures all services are ready before any user message
arrives.

\begin{lstlisting}[language=Python, caption={Startup sequence (simplified from main.py)}]
async def setup_hook(self):
    # 1. Relational database
    self.db = Database(config.database_path)
    self.db.connect()
    self.db.init_tables()

    # 2. Knowledge base (loads .md and .yml files into memory)
    self.kb = KnowledgeBase(data_dir=config.data_dir)
    self.kb.load()

    # 3. Fuzzy search over loaded KB
    self.search_svc = SearchService(self.kb)

    # 4. Rate limiter (in-memory, resets on restart)
    self.rate_limiter = RateLimiter(max_calls=5, period_seconds=60)

    # 5. LLM client (only if OLLAMA_ENABLED=true)
    self.ollama = OllamaClient(...)

    # 6. Conversation manager (in-memory per-user sessions)
    self.conversation_manager = ConversationManager(...)

    # 7. Embedding service (nomic-embed-text via Ollama)
    self.embedder = EmbeddingService(...)

    # 8. ChromaDB vector store (reads from ./chroma_db/)
    self.vector_store = VectorStore(db_path=config.chroma_db_path)

    # 9. Retriever (ties embedder + vector store together)
    self.retriever = Retriever(embedder=self.embedder,
                               vector_store=self.vector_store)

    # 10. Intent detector
    self.intent_detector = IntentDetector()

    # 11. Load all cogs (commands)
    for ext in EXTENSIONS:
        await self.load_extension(ext)
\end{lstlisting}

All services are attached to the \texttt{bot} instance. Every cog accesses them
via \texttt{self.bot.db}, \texttt{self.bot.retriever}, etc.

\subsection{Message Routing Flow}

When a student sends a message in \#ask-gsa (or DMs the bot), the
\texttt{ChatCog.on\_message} handler processes it through a series of gates:

\begin{lstlisting}[language=bash, caption={Message routing gates (chat.py)}]
Gate 1: Is this the bot's own message?      -> ignore
Gate 2: Is this a slash command?            -> ignore (handled by cog)
Gate 3: Does message mention another user?  -> ignore (avoid cross-talk)
Gate 4: Is it the ask-gsa channel or DM?   -> proceed | ignore
Gate 5: IntentDetector classifies intent
  - greeting    -> warm reply, no RAG
  - thanks      -> acknowledgement, no RAG
  - food        -> retrieve events with food tag
  - social      -> retrieve social events
  - clear       -> wipe session history
  - help        -> show command list
  - question/statement -> full RAG pipeline
Gate 6: Rate limiter check (5 req/60s)      -> throttle if exceeded
Gate 7: RAG pipeline                        -> generate answer
\end{lstlisting}

% =============================================================================
\section{The RAG Pipeline — End to End}
% =============================================================================

Retrieval-Augmented Generation (RAG) is the core intelligence of the bot.
Instead of fine-tuning a model on GSA data (expensive, hard to update), RAG
retrieves relevant documents at query time and injects them into the LLM prompt.
The LLM is explicitly instructed to answer \emph{only} from the provided
documents, preventing hallucination.

\subsection{Offline Phase: Building the Index}

This runs once (or after any KB file edit) via
\texttt{python scripts/build\_index.py --reset}.

\begin{figure}[H]
\centering
\begin{mdframed}
\begin{lstlisting}[language=bash, numbers=none]
  Knowledge Base Files
  (.md, .yml in bot/data/)
         |
         v
  [DocumentChunker]
  Splits each file into <= 350-token chunks
  Each chunk: text + source_file + section_title + metadata
         |
         v
  [EmbeddingService]
  Calls Ollama /api/embed with "search_document: {text}"
  Model: nomic-embed-text  |  Output: 768-dimensional float vector
         |
         v
  [VectorStore (ChromaDB)]
  Stores: chunk_id, text, embedding, metadata
  Location: ./chroma_db/   (persistent, survives restart)
  143 chunks total from 10 source files
\end{lstlisting}
\end{mdframed}
\caption{Offline indexing pipeline}
\end{figure}

\subsection{Online Phase: Answering a Question}

This runs for every student message that reaches the RAG path.

\begin{figure}[H]
\centering
\begin{mdframed}
\begin{lstlisting}[language=bash, numbers=none]
  Student message: "How many speakers at MMI 2026 from Rutgers?"
         |
         v
  [OllamaClient.expand_query()]          (optional)
  Short/vague queries are rewritten:
  "MMI" -> "What is the MMI Workshop Series at NJIT?"
         |
         v
  [Retriever._build_search_query()]
  Cleans mentions/channels, optionally appends
  topic keywords from conversation history
         |
         v
  PARALLEL:
  [EmbeddingService.embed_query()]        [BM25Index.search()]
  Semantic vector search                  Lexical exact-term search
  "search_query: " prefix                 tokenize + BM25Okapi scores
  top-20 by cosine similarity             top-20 by BM25 score
         |                                       |
         +---------------+   +-------------------+
                         |   |
                         v   v
  [Retriever._reciprocal_rank_fusion()]
  Merge both ranked lists: score = sum(1 / (60 + rank))
  Normalize fused scores to [0.40, 0.95] -> stored as rrf_score
  BM25-only chunks (unknown terms, sparse matches) are included
  Result: up to 40 unique candidates
         |
         v
  [Retriever._rerank()]
  base_score   = rrf_score (not raw cosine -- fair for BM25 chunks)
  common_bonus = word-boundary hits * 0.05 (max common contribution)
  proper_bonus = capitalized non-stopword hits * 0.15 (3x weight:
                 names like "Singh", "MARCuS" are key disambiguators)
  source_bonus = faq:0.05, policy:0.03, event:0.02
  title_bonus  = 0.10 if query word in section title
  Keeps top-7 by relevance_score
         |
         v
  [Retriever._include_siblings()]
  For any chunk tagged with qa_base_id:
  fetch all continuation parts of split Q&A answers
         |
         v
  [OllamaClient.generate_answer()]
  SYSTEM: GSA Gateway persona + 11 rules (incl. exhaustive enumeration)
          + conversation history
  USER:   [Document 1...Document N] + student question
  Calls Ollama /api/generate  (llama3.1:8b, temp=0.3)
         |
         v
  Discord reply with source attribution
  "According to Document 2: MMI Workshop Series..."
\end{lstlisting}
\end{mdframed}
\caption{Online hybrid retrieval and generation pipeline (as of June 2026)}\end{figure}

\subsection{The Embedding Asymmetry}

A subtle but important design decision: documents and queries are embedded with
different prefixes. This is specific to \texttt{nomic-embed-text}:

\begin{center}
\begin{tabular}{ll}
\toprule
\textbf{Case} & \textbf{Prefix added before embedding} \\
\midrule
Indexing a document chunk & \texttt{"search\_document: \{text\}"} \\
Embedding a user query & \texttt{"search\_query: \{text\}"} \\
\bottomrule
\end{tabular}
\end{center}

This is called \emph{asymmetric dense retrieval}. The model was trained to align
query vectors with document vectors even when the phrasing differs. Without the
correct prefix, retrieval quality degrades noticeably.

\subsection{Chunking Strategy}

The \texttt{DocumentChunker} uses different strategies per file type:

\subsubsection{Markdown FAQ Files (\texttt{chunk\_markdown\_faq})}

Splits on \texttt{\#\#} headings. Each Q\&A pair becomes one chunk. If the full
\texttt{"Question: ...\textbackslash nAnswer: ..."} text exceeds 350 tokens, the
answer is split using \texttt{split\_text\_by\_tokens()}, which:
\begin{itemize}[leftmargin=2em]
  \item Splits text into sentences (on \texttt{.!?} boundaries and newlines)
  \item Accumulates sentences until the 350-token limit is approached
  \item Carries a 50-token overlap into the next chunk for context continuity
  \item Tags multi-part chunks with \texttt{qa\_base\_id} and \texttt{chunk\_index}
        metadata so the retriever can always include all parts of a split answer
\end{itemize}

\subsubsection{Markdown Policy Files (\texttt{chunk\_markdown\_policy})}

Splits on \texttt{\#\#} headings. Each section becomes one chunk. If a section
exceeds 350 tokens, it is sub-chunked with the same sentence-aware splitter.
Section titles are prepended to each sub-chunk for context.

\subsubsection{YAML Events (\texttt{chunk\_yaml\_events})}

One chunk per event, formatted as structured text:
\begin{lstlisting}[numbers=none]
Event: MMI Workshop 2026
Date: 2026-05-21
Location: NJIT GITC 1100
Description: ...
\end{lstlisting}

\subsubsection{YAML Contacts (\texttt{chunk\_yaml\_contacts})}

One chunk per officer entry, one chunk per campus office entry. Each contains
role, name, email, responsibilities.

\subsubsection{YAML Resources (\texttt{chunk\_yaml\_resources})}

One chunk per resource item across all categories. Each chunk includes the
category name for context.

\subsection{Conversation Memory}

\texttt{ConversationManager} maintains per-user session history entirely
\textbf{in memory} — there is no database table for conversations. This is
intentional: session data is private by design and does not need to persist
across restarts.

\begin{center}
\begin{tabular}{ll}
\toprule
\textbf{Setting} & \textbf{Value (configurable in .env)} \\
\midrule
Session timeout & 60 minutes of inactivity \\
Max turns kept & 5 user+assistant pairs (10 total turns) \\
Cleanup interval & Every 5 minutes (background async task) \\
Storage & Python dict in-memory, never on disk \\
\bottomrule
\end{tabular}
\end{center}

When conversation history exists, two things change:
\begin{enumerate}[leftmargin=2em]
  \item The retriever appends up to 5 topic keywords from recent assistant
        responses to the search query, improving follow-up retrieval.
  \item The LLM system prompt is extended with the last 5 turns so it can
        resolve references like ``what about step 2?'' or ``how much is that?''
\end{enumerate}

\subsection{Query Expansion}

When a student types something very short or ambiguous (e.g., ``MMI'',
``funding'', ``travel''), a preliminary call to \texttt{expand\_query()} rewrites
it to a full question before retrieval. This prevents poor embedding alignment
from a single token query.

\begin{lstlisting}[numbers=none]
Input:  "MMI"
Output: "What is the MMI Workshop Series at NJIT?"

Input:  "travel"
Output: "How can I apply for a GSA travel award?"
\end{lstlisting}

The expansion uses \texttt{llama3.1:8b} with temperature 0.1 and \texttt{num\_predict=60},
making it fast. If Ollama is unavailable, the original query is used as-is.

% =============================================================================
\section{Vector Database: ChromaDB}
% =============================================================================

\subsection{What is a Vector Database?}

A vector database stores high-dimensional float vectors (embeddings) and
retrieves the nearest neighbours by distance metric. In this system:

\begin{itemize}[leftmargin=2em]
  \item Every KB chunk is converted to a 768-dimensional vector by
        \texttt{nomic-embed-text}
  \item At query time, the student's question is also converted to a vector
  \item The database returns the chunks whose vectors are closest to the query
        vector (cosine similarity)
  \item Semantic similarity is captured, not just keyword matching
\end{itemize}

\subsection{Why ChromaDB?}

\begin{center}
\begin{longtable}{p{3cm}p{10cm}}
\toprule
\textbf{Reason} & \textbf{Detail} \\
\midrule
No server required & ChromaDB runs in-process. You do not need to start a
  separate database server. For a bot running on a single machine, this
  eliminates an entire infrastructure dependency. \\
\midrule
Persistent by default & \texttt{chromadb.PersistentClient(path="./chroma\_db")}
  writes to disk automatically. The index survives restarts without any extra
  configuration. \\
\midrule
Python-native API & Direct Python calls with no query language. Adding chunks,
  querying, filtering, and resetting are all single function calls. \\
\midrule
Metadata filtering & Supports \texttt{where=\{"source\_type": "faq"\}} style
  filters, which we use to restrict food queries to event chunks only. \\
\midrule
Small-to-medium scale & For 143 chunks (our current size), ChromaDB uses HNSW
  indexing with cosine space, which is extremely fast. Even at 10,000 chunks it
  remains suitable. \\
\midrule
Open source, free & MIT license. No usage limits, no API keys. \\
\bottomrule
\end{longtable}
\end{center}

\subsection{ChromaDB in Our Codebase}

The \texttt{VectorStore} class (\texttt{bot/services/vector\_store.py}) wraps
all ChromaDB interaction:

\begin{lstlisting}[language=Python, caption={Key VectorStore methods}]
class VectorStore:
    def add_chunks(chunks, embeddings)
        # Bulk add during index build

    def query(query_embedding, n_results=20, source_type_filter=None)
        # Vector similarity search, returns list[dict] with similarity scores

    def get_all_chunks()
        # Return every chunk -- used to build the BM25 index at startup

    def get_sibling_chunks(qa_base_id)
        # Fetch all continuation parts of a split Q&A answer

    def get_chunk_count()
        # Total chunks in the collection

    def reset()
        # Delete and recreate collection (used by build_index.py --reset)
\end{lstlisting}

ChromaDB stores three things per chunk:
\begin{enumerate}[leftmargin=2em]
  \item \textbf{ID} — unique string identifier (e.g., \texttt{mmi\_workshop\_faq\_7\_0\_1})
  \item \textbf{Document} — the raw chunk text
  \item \textbf{Metadata} — structured fields: \texttt{source\_file},
        \texttt{source\_type}, \texttt{section\_title}, \texttt{token\_count},
        and optionally \texttt{qa\_base\_id} + \texttt{chunk\_index} for
        multi-part answers
\end{enumerate}

\subsection{Alternatives to ChromaDB}

Understanding why we chose ChromaDB requires understanding the alternatives.

\begin{center}
\begin{longtable}{p{2.5cm}p{4cm}p{4cm}p{2.5cm}}
\toprule
\textbf{Database} & \textbf{Strengths} & \textbf{Weaknesses} & \textbf{Best for} \\
\midrule
\textbf{ChromaDB}
  & No server, persistent, Python-native, free
  & Not suitable for multi-node or billion-scale
  & Our use case \\
\midrule
\textbf{FAISS} (Meta)
  & Extremely fast, industry-proven, supports GPU indexing
  & In-memory only by default (no built-in persistence), no metadata filtering, lower-level API
  & High-performance search where you manage persistence yourself \\
\midrule
\textbf{Qdrant}
  & High performance, rich filtering, Rust internals, can run as embedded or server
  & Slightly more complex setup than ChromaDB
  & Production systems needing advanced filters \\
\midrule
\textbf{Weaviate}
  & GraphQL API, multi-modal, schema-driven
  & Requires running a separate server, heavier infrastructure
  & Enterprise, multi-modal (image+text) retrieval \\
\midrule
\textbf{Pinecone}
  & Fully managed cloud, serverless option, scales automatically
  & Paid after free tier, data leaves your infrastructure
  & Teams that do not want to manage infrastructure \\
\midrule
\textbf{Milvus}
  & Billion-scale, distributed, battle-tested at Zilliz
  & Heavy infrastructure (requires etcd, MinIO, Pulsar)
  & Large-scale enterprise with dedicated DevOps \\
\midrule
\textbf{pgvector}
  & Lives inside PostgreSQL, SQL-native, easy to join with relational data
  & Requires PostgreSQL, slower than dedicated vector DBs at scale
  & Teams already running PostgreSQL \\
\bottomrule
\end{longtable}
\end{center}

\begin{infobox}[When you would migrate away from ChromaDB]
If the knowledge base grows beyond $\sim$50,000 chunks, or if you need to
deploy on multiple machines, Qdrant (embedded mode) or pgvector are the natural
next steps. Both can be swapped in by replacing only \texttt{vector\_store.py} —
no other code changes needed.
\end{infobox}

\subsection{The Distance Metric: Cosine Similarity}

ChromaDB is configured with \texttt{"hnsw:space": "cosine"}. Cosine similarity
measures the angle between two vectors, ignoring magnitude:

\[
\text{cosine\_similarity}(A, B) = \frac{A \cdot B}{\|A\| \cdot \|B\|}
\]

A score of 1.0 means the vectors point in exactly the same direction (identical
semantic meaning). A score of 0.0 means perpendicular (unrelated). ChromaDB
returns distances (\texttt{1.0 - similarity}), which we convert back:
\texttt{similarity = 1.0 - distance}.

Our minimum similarity threshold is \textbf{0.3}. Chunks below this are
discarded before reranking.

% =============================================================================
\section{The Embedding Model: nomic-embed-text}
% =============================================================================

\texttt{nomic-embed-text} is an open-source text embedding model developed by
Nomic AI. It produces 768-dimensional vectors and is specifically designed for
retrieval tasks.

\subsection{Why This Model?}

\begin{itemize}[leftmargin=2em]
  \item Runs locally via Ollama — no API key, no network call, no cost
  \item Strong performance on standard retrieval benchmarks (BEIR suite)
  \item Supports asymmetric retrieval via task prefixes
        (\texttt{search\_query:} vs \texttt{search\_document:})
  \item Fast enough to embed 143 chunks in under 1 second on a CPU
\end{itemize}

\subsection{Alternatives}

\begin{center}
\begin{tabular}{llll}
\toprule
\textbf{Model} & \textbf{Dimensions} & \textbf{Local?} & \textbf{Cost} \\
\midrule
nomic-embed-text (ours) & 768 & Yes (Ollama) & Free \\
text-embedding-3-small (OpenAI) & 1536 & No & \$0.02/1M tokens \\
text-embedding-3-large (OpenAI) & 3072 & No & \$0.13/1M tokens \\
mxbai-embed-large & 1024 & Yes (Ollama) & Free \\
all-MiniLM-L6-v2 (sentence-transformers) & 384 & Yes & Free \\
\bottomrule
\end{tabular}
\end{center}

To swap the embedding model, change \texttt{EMBEDDING\_MODEL} in \texttt{.env}
and rebuild the index. The vector store must be fully rebuilt because vector
dimensions and semantic spaces differ between models.

% =============================================================================
\section{The LLM: Ollama + llama3.1:8b}
% =============================================================================

\subsection{What Ollama Is}

Ollama is a local model server that downloads and serves open-weight language
models. The bot talks to it via HTTP on \texttt{localhost:11434}.

\begin{lstlisting}[language=bash, caption={Ollama endpoints used}]
POST /api/generate   # Text generation (question answering, query expansion)
POST /api/embed      # Text embedding (used by EmbeddingService)
\end{lstlisting}

\subsection{The System Prompt Design}

\texttt{OllamaClient} constructs a two-part prompt for every RAG answer:

\textbf{System prompt} (always present):
\begin{itemize}[leftmargin=2em]
  \item Defines the GSA Gateway persona
  \item Lists 10 explicit rules (never hallucinate, cite documents, format for Discord, etc.)
  \item Appends conversation history when a session exists
\end{itemize}

\textbf{User prompt} (built per query):
\begin{itemize}[leftmargin=2em]
  \item Injects retrieved chunks as numbered documents with source names and relevance scores
  \item Appends the student's question
  \item Instructs the model to answer only from the provided documents
\end{itemize}

Generation parameters: \texttt{temperature=0.3} (low randomness for factual
answers), \texttt{top\_p=0.9}, \texttt{num\_predict=512} (max response tokens).

\subsection{Graceful Degradation}

If Ollama is unavailable or times out (60-second default), the bot falls back to
the \texttt{SearchService} fuzzy-search answer. If fuzzy search also yields no
result above 60\% confidence, the bot returns a polite referral message to a
GSA officer. The system never silently fails.

% =============================================================================
\section{The Relational Database: SQLite}
% =============================================================================

SQLite stores everything that needs to persist across restarts and that has
structured, queryable form.

\subsection{Tables}

\begin{center}
\begin{tabular}{ll}
\toprule
\textbf{Table} & \textbf{Contents} \\
\midrule
\texttt{feedback} & Anonymous student feedback submissions \\
\texttt{initiatives} & Initiative proposals (5-field Modal form) \\
\texttt{events} & Events added via \texttt{/admin\_add\_event} \\
\texttt{message\_stats} & Per-channel message counts for weekly summaries \\
\bottomrule
\end{tabular}
\end{center}

\subsection{Privacy Design}

User IDs are never stored in plain text. Before any database write,
\texttt{hash\_user\_id(user\_id)} applies a SHA-256 hash (with a secret salt
from \texttt{.env}). This means GSA officers can see aggregate statistics
(``12 feedback submissions this week'') but cannot identify which student
submitted which entry.

% =============================================================================
\section{Fuzzy Search: rapidfuzz}
% =============================================================================

\texttt{SearchService} (\texttt{bot/services/search.py}) provides keyword-based
approximate matching as a fallback when RAG is unavailable or as a complement to
it in the \texttt{/ask} command.

\texttt{rapidfuzz} uses the WRatio algorithm (Weighted Ratio), which combines
multiple string similarity measures and returns a score from 0–100. The
confidence threshold is \textbf{60.0}: matches below this are discarded and a
fallback message is shown.

\begin{lstlisting}[language=Python, caption={How the threshold works}]
# In search.py
result = process.extractOne(query, candidates, scorer=fuzz.WRatio)
if result and result[1] >= 60.0:
    return best_match
else:
    return None  # triggers fallback message
\end{lstlisting}

% =============================================================================
\section{Configuration System}
% =============================================================================

All configuration lives in a \texttt{.env} file (never committed to git). The
\texttt{Config} dataclass in \texttt{bot/config.py} reads these at startup and
provides typed access throughout the codebase.

\begin{lstlisting}[language=bash, caption={Key .env variables}]
# Required
DISCORD_TOKEN=your_token_here
DISCORD_GUILD_ID=your_guild_id

# Ollama
OLLAMA_ENABLED=true
OLLAMA_MODEL=llama3.1:8b
OLLAMA_URL=http://localhost:11434
OLLAMA_TIMEOUT=60
EMBEDDING_MODEL=nomic-embed-text

# Bot behavior
ADMIN_ROLE_NAME=GSA Officer
ALLOWED_CHANNELS=ask-gsa,gsa-bot-testing
ASK_GSA_CHANNEL=ask-gsa
BOT_PREFIX=gsa

# RAG / vector store
CHROMA_DB_PATH=./chroma_db
CONVERSATION_TIMEOUT_MINUTES=60
CONVERSATION_MAX_TURNS=5

# Database and privacy
DATABASE_PATH=./gsa_gateway.db
USER_ID_SALT=your_random_salt_here

# Channel names (must match Discord channel names exactly)
CHANNEL_ANNOUNCEMENTS=gsa-announcements
CHANNEL_EVENTS=gsa-events
CHANNEL_FOOD=gsa-food
MATHCAFE_CHANNEL=gsa-mathcafe
\end{lstlisting}

% =============================================================================
\section{Scripts}
% =============================================================================

\begin{center}
\begin{longtable}{lp{9cm}}
\toprule
\textbf{Script} & \textbf{Purpose} \\
\midrule
\texttt{scripts/setup.sh} & Full first-time setup: creates venv, installs dependencies,
  initialises SQLite tables, runs tests. Run once on a new machine. \\
\texttt{scripts/build\_index.py} & Rebuilds the ChromaDB vector index from all
  \texttt{bot/data/} files. Run after any KB file edit.
  \texttt{--reset} flag forces full rebuild without prompting. \\
\texttt{scripts/run\_bot.sh} & Starts the bot process with log output to
  \texttt{gsa\_gateway.log}. \\
\texttt{scripts/restart.sh} & Stops any running bot process, optionally restarts
  Ollama, and starts the bot fresh. Safe to run during active sessions. \\
\texttt{scripts/health\_check.sh} & Checks whether the bot process is running,
  how long it has been up, memory usage, and MathCafe last-post status. \\
\texttt{scripts/export\_events\_json.py} & Syncs \texttt{events.yml} to
  \texttt{website/data/events.json} for the static website. Run after
  adding events. \\
\texttt{scripts/export\_weekly\_summary.py} & Generates and prints the 7-day
  admin summary to stdout or a file. \\
\texttt{scripts/init\_db.py} & Creates SQLite tables. Safe to re-run
  (idempotent). \\
\bottomrule
\end{longtable}
\end{center}

% =============================================================================
\section{The Static Website}
% =============================================================================

The website at \texttt{website/} is pure HTML, CSS, and JavaScript with no
build step and no framework. It is designed to be hosted on GitHub Pages.

\begin{center}
\begin{tabular}{ll}
\toprule
\textbf{File} & \textbf{Content} \\
\midrule
\texttt{index.html} & Hero section, feature cards, slash command table \\
\texttt{about.html} & GSA mission, officer bios \\
\texttt{events.html} & Loads and renders \texttt{data/events.json} via \texttt{fetch()} \\
\texttt{initiatives.html} & Call to action for using \texttt{/initiative} in Discord \\
\texttt{resources.html} & Static resource listings by category \\
\texttt{contact.html} & Officer and campus office directory \\
\texttt{assets/style.css} & NJIT red \texttt{\#CC0000} + dark gray, mobile-first, no frameworks \\
\texttt{assets/app.js} & Nav toggle, events loader, HTML escaping for XSS safety \\
\texttt{data/events.json} & Auto-exported from \texttt{events.yml} by the export script \\
\bottomrule
\end{tabular}
\end{center}

\begin{warningbox}[Keeping events in sync]
The website reads \texttt{events.json}, not \texttt{events.yml}. After editing
\texttt{events.yml}, you must run \texttt{python scripts/export\_events\_json.py}
and commit the updated \texttt{events.json}. If you forget, the website will
show stale events.
\end{warningbox}

% =============================================================================
\section{MathCafe Feature}
% =============================================================================

\texttt{MathCafeService} (\texttt{bot/services/mathcafe.py}) posts a daily
mathematics fact to the \texttt{\#gsa-mathcafe} Discord channel. Facts are drawn
from a queue in \texttt{bot/data/mathcafe/facts.yml}.

Special behavior: on \textbf{May 12th} (Maryam Mirzakhani's birthday), the bot
posts a dedicated tribute fact using \texttt{bot/data/mathcafe/images/Maryam.png}
and a biographical message regardless of the regular queue. It tracks
\texttt{last\_posted\_year} to ensure this fires exactly once per year.

% =============================================================================
\section{Security and Privacy Design}
% =============================================================================

\begin{enumerate}[leftmargin=2em]
  \item \textbf{User IDs are hashed.} No Discord user ID is ever stored in
        plain text. SHA-256 with a secret salt is applied in
        \texttt{database.py:hash\_user\_id()}.
  \item \textbf{Admin commands are ephemeral.} All responses in \texttt{admin.py}
        use \texttt{ephemeral=True}, meaning only the requesting officer sees them.
  \item \textbf{Rate limiting} prevents abuse: 5 RAG requests per user per 60 seconds.
  \item \textbf{Channel allowlist} restricts bot responses to named channels only.
  \item \textbf{LLM grounding.} The system prompt explicitly forbids the LLM from
        answering outside the provided context, preventing prompt injection from
        fabricating GSA policies.
  \item \textbf{Admin role gate.} Admin commands check the \texttt{ADMIN\_ROLE\_NAME}
        Discord role before executing.
  \item \textbf{Message gate 3.} Messages that mention another user (via
        \texttt{message.mentions} or raw \texttt{<@...>}) are ignored when the
        bot was not directly mentioned, preventing the bot from joining
        cross-user conversations uninvited.
\end{enumerate}

% =============================================================================
\section{Running the System: Quick Reference}
% =============================================================================

\subsection{First-Time Setup}

\begin{lstlisting}[language=bash]
# 1. Clone and enter the repo
git clone <repo-url> gsa-gateway
cd gsa-gateway

# 2. Run the setup script (creates venv, installs deps, inits DB)
bash scripts/setup.sh

# 3. Copy and fill in the environment file
cp .env.example .env
# Edit .env with your DISCORD_TOKEN and settings

# 4. Pull required Ollama models
ollama pull llama3.1:8b
ollama pull nomic-embed-text

# 5. Build the vector index
python3 scripts/build_index.py --reset

# 6. Start the bot
bash scripts/run_bot.sh
\end{lstlisting}

\subsection{After Editing Knowledge Base Files}

\begin{lstlisting}[language=bash]
# Rebuild vector index (required after any .md or .yml edit in bot/data/)
python3 scripts/build_index.py --reset

# If you added events, also sync the website
python3 scripts/export_events_json.py

# Restart the bot to reload the knowledge base into memory
bash scripts/restart.sh
\end{lstlisting}

\subsection{Health Check}

\begin{lstlisting}[language=bash]
bash scripts/health_check.sh          # Process status, uptime, memory
tail -f gsa_gateway.log               # Live log stream
# In Discord (admin only):
/admin_stats                          # RAG chunk count, active sessions
\end{lstlisting}

% =============================================================================
\section{Testing}
% =============================================================================

Tests live in \texttt{bot/tests/} and use \texttt{pytest}. They operate on
in-memory fixtures and do not require Ollama or Discord to be running.

\begin{center}
\begin{tabular}{ll}
\toprule
\textbf{File} & \textbf{What It Tests} \\
\midrule
\texttt{conftest.py} & Fixtures: in-memory SQLite, KB, SearchService \\
\texttt{test\_search.py} & Fuzzy search hits, misses, confidence thresholds \\
\texttt{test\_database.py} & CRUD operations, privacy hashing, aggregate stats \\
\texttt{test\_commands.py} & Admin role check, rate limiter, channel allowlist \\
\bottomrule
\end{tabular}
\end{center}

\begin{lstlisting}[language=bash]
# Run all tests
pytest bot/tests/ -v

# Run a specific test file
pytest bot/tests/test_search.py -v
\end{lstlisting}

% =============================================================================
\section{Key Invariants for Contributors}
% =============================================================================

These are rules the system depends on. Breaking them causes silent failures,
not errors.

\begin{infobox}[Must-know rules]
\begin{enumerate}[leftmargin=1.5em]
  \item \textbf{Never store raw user IDs.} Always call \texttt{hash\_user\_id()}
        before writing to the database.
  \item \textbf{Admin commands must use \texttt{ephemeral=True}.} Officer-only
        data should never be visible to the whole channel.
  \item \textbf{Rebuild the index after any KB edit.} The ChromaDB store does
        not watch files; it must be explicitly rebuilt.
  \item \textbf{The embedding prefix matters.} Documents use
        \texttt{"search\_document: "}, queries use \texttt{"search\_query: "}.
        Swapping them degrades retrieval silently.
  \item \textbf{Conversation sessions are in-memory only.} They reset on restart.
        This is a feature, not a bug.
  \item \textbf{The LLM never answers without context.} If \texttt{chunks} is
        empty, \texttt{generate\_answer()} returns \texttt{None} and the bot
        falls back to fuzzy search.
  \item \textbf{All config comes from \texttt{.env}.} Never hardcode tokens,
        paths, or channel IDs in Python files.
\end{enumerate}
\end{infobox}

% =============================================================================
\section{Architecture Decision Log}
% =============================================================================

This section explains \emph{why} specific choices were made, so future
contributors do not reverse them accidentally.

\begin{enumerate}[leftmargin=2em]
  \item \textbf{Why SQLite and not PostgreSQL?} The bot runs on a single
        machine with one writer at a time (the bot process). SQLite is
        sufficient, requires no server, and makes deployment trivial. If
        the bot is ever deployed across multiple servers, PostgreSQL would
        be the migration target.

  \item \textbf{Why Ollama and not OpenAI API?} Two reasons: cost (the
        bot serves students continuously; API calls accumulate) and privacy
        (student questions about visa status, mental health resources, and
        finances should not leave the server). Local inference with
        \texttt{llama3.1:8b} covers the quality bar for factual Q\&A.

  \item \textbf{Why ChromaDB and not FAISS?} FAISS does not persist
        to disk natively and has no metadata filtering. ChromaDB does both
        out of the box. For 143 chunks, performance difference is
        immeasurable.

  \item \textbf{Why \texttt{nomic-embed-text} and not a sentence-transformer?}
        It runs via the same Ollama endpoint already in use for generation,
        keeping infrastructure uniform. It also natively supports asymmetric
        retrieval prefixes, which measurably improves recall for short
        student queries against longer document chunks.

  \item \textbf{Why 350-token chunk limit?} Empirically chosen to balance
        two competing forces: embedding quality degrades for very long texts
        (over $\sim$500 tokens), and retrieval misses context if chunks are
        too short (under $\sim$100 tokens). 350 tokens keeps the typical
        Q\&A answer together while fitting within the model's sweet spot.

  \item \textbf{Why sibling chunk expansion (added June 2026)?} Long Q\&A
        answers that exceed 350 tokens are split into continuation chunks.
        The reranker would sometimes drop the second part because the first
        part scored higher (more keyword hits). The sibling system ensures
        that whenever any part of a split answer is retrieved, all parts
        are included. This fixed a bug where Prof.\ Enrique Dunn (Day 2
        schedule) was invisible to queries about Stevens speakers.

  \item \textbf{Why hybrid BM25 + vector retrieval (added June 2026)?}
        Pure vector search has two systematic failure modes. First, the
        embedding model has no representation for novel terms (acronyms,
        technical strings like ``InstructEx'', proper names). These terms
        score identically across all chunks by cosine similarity, so the
        relevant chunk is never retrieved. Second, when many chunks have
        near-identical cosine scores, the wrong chunk surfaces. BM25 lexical
        search solves the first problem via exact term matching; the proper
        noun bonus (3$\times$ weight) in the reranker solves the second.
        Reciprocal Rank Fusion merges both lists with a normalized score
        used as the reranker base, so BM25-only chunks are not penalized
        by a placeholder similarity value. Tested against 10 hard questions;
        all answered correctly after implementation.

  \item \textbf{Why BM25Okapi and not BM25Plus or BM25L?} BM25Okapi is the
        standard formulation used in most IR literature and is the default
        in \texttt{rank\_bm25}. For the KB size (165 chunks), the length
        normalization differences between variants are negligible. BM25Okapi
        was sufficient once the RRF score normalization fixed the base-score
        collapse for BM25-only results. Migrate to BM25Plus if the KB grows
        significantly and long chunks are consistently under-ranked.
\end{enumerate}

% =============================================================================
\section{Glossary}
% =============================================================================

\begin{description}[leftmargin=3em, style=nextline]
  \item[RAG (Retrieval-Augmented Generation)] Architecture where a retrieval
    step fetches relevant documents before generation. The LLM is grounded in
    retrieved facts rather than relying solely on training data.
  \item[Embedding] A high-dimensional float vector that represents the semantic
    meaning of a piece of text. Similar texts produce similar vectors.
  \item[Cosine similarity] Distance metric between two vectors measured as the
    cosine of the angle between them. Value range: $[0, 1]$ after normalization.
  \item[Chunk] A fragment of a knowledge base document, sized to fit within a
    token limit, that becomes one entry in the vector store.
  \item[HNSW] Hierarchical Navigable Small World — the approximate nearest
    neighbour algorithm used by ChromaDB for fast vector search.
  \item[Cog] discord.py term for a module that groups related commands and
    event listeners. Loaded and unloaded as a unit.
  \item[Ephemeral message] A Discord message visible only to the user who
    triggered the command, not to the rest of the channel.
  \item[Intent detection] Classifying the purpose of an incoming message
    (greeting, question, food query, etc.) to route it to the right handler
    without involving the LLM for simple cases.
  \item[Reranking] A second scoring pass after initial retrieval. Combines
    cosine similarity with keyword density, source type, and section title
    matches to improve the ordering of retrieved chunks.
  \item[qa\_base\_id] Metadata tag added to continuation chunks of a multi-part
    Q\&A answer. Enables the retriever to fetch all sibling parts when any
    one part is retrieved.
  \item[Asymmetric retrieval] Using different prefix strings when embedding
    queries vs.\ documents. Exploits model training to better align short
    questions with longer document passages.
  \item[BM25 (Best Match 25)] Probabilistic lexical retrieval algorithm that
    scores chunks by term frequency and inverse document frequency. Handles
    exact/sparse term matching that semantic embeddings cannot. Used alongside
    vector search in hybrid retrieval.
  \item[Reciprocal Rank Fusion (RRF)] A method for merging multiple ranked
    lists. Each item scores $\sum 1/(k + \text{rank})$. Items appearing in
    both vector and BM25 results receive combined scores; BM25-only items
    (exact term matches) are still included. $k{=}60$ is the standard constant.
  \item[Hybrid retrieval] Running vector search and BM25 in parallel and
    fusing the results. Fixes two failure classes: novel/unknown terms and
    near-equal cosine similarity ties.
\end{description}

% =============================================================================
\vspace{2em}
\begin{center}
\rule{0.5\textwidth}{0.4pt}\\[0.5em]
{\small GSA Gateway — NJIT Graduate Student Association\\
Contact: \href{mailto:md724@njit.edu}{md724@njit.edu} $\cdot$
Discord: \texttt{discord.gg/a4mvbEmSAq} $\cdot$
June 2026}
\end{center}

\end{document}
